% =====================================================================
% CT v1.3 §3.5.1 — Okamoto-rename of the V_quad parameter point
% (η, θ) → (α, β)
%
% STATUS: AUTHORED 2026-05-08 by SIARC agent (Copilot Researcher tier)
%         per relay 105 envelope. Closes 069r2 QB.3 = Y_RENAME_REQUIRED.
%
% INSERT POSITION: pcf-research/channel/cc_pipeline_v13_2026-05-02/
%   channel_theory_outline.tex
%   immediately after line 891 (end of §3.5 prose), before the
%   "% =====" divider on line 893 and the §3.6 \subsection on line 894.
%
% CITATIONS USED: only existing bibkeys
%   {iwasaki1991from, ohyama2006studies, noumiyamada1998painleve,
%    jimbomiwa1981monodromy, siarc_cc_pipeline_g}.
% NEW BIBKEYS RECOMMENDED (not introduced in this insert; documented
% as UF-105-NEW-BIBKEY in unexpected_finds.json for operator splice):
%   {okamoto1987painleveIII, kajiwaranoumiyamada2017,
%    forresterwitte2002}.
% Inline-prose attribution is used for these three primary sources to
% preserve pdflatex compile cleanliness with no undefined references.
% =====================================================================

\subsubsection*{Okamoto-rename of the $\Vquad$ parameter point}
\label{sssec:okamoto-rename}

The $\Vquad \rightsquigarrow P_{III}(D_6)$ reduction of
\Cref{ex:cc} and \Cref{sec:position}(b) is anchored at the
parameter point $(1/6,\,0,\,0,\,-1/2)$ in $\mathbb{Q}^4$, with
the symbol convention $(\alpha_\infty, \alpha_0, \beta_\infty,
\beta_0)$. The convention itself has, until v1.3, been used
without an explicit introduction. We supply that introduction
here, and record the offset of $-1/3$ that the $\Vquad$ image
carries against the canonical Okamoto null-sum constraint.

\medskip
\noindent\textbf{Background: three coexisting Painlev\'e
parameter conventions.}\hspace{0.5em}
The standard Hamiltonian theory of $P_{III}$ (Okamoto,
{\em Funkcial.\ Ekvac.}\ {\bf 30} (1987), 305--332, \S 1.1)
writes the Hamiltonian
\[
  H_{III} \;=\; \tfrac{1}{t}\bigl[2 q^{2} p^{2} -
       \{2 \eta_\infty\, t q^{2} + (2 \theta_0 + 1) q
       - 2 \eta_0\, t\}p
       + \eta_\infty(\theta_0 + \theta_\infty)\, t q\bigr]
\]
in the four parameters $(\eta_\infty, \eta_0, \theta_\infty,
\theta_0)$, with the WLOG normalisation
$\eta_\infty = \eta_0 = 1$ adopted by Okamoto immediately
after this display (loc.\ cit., eq.\ (0.5)). The
Sakai-classification treatment of the
$((2 A_1)^{(1)} / D_6^{(1)})$ surface (Kajiwara, Noumi, Yamada,
{\em J.\ Phys.\ A: Math.\ Theor.}\ {\bf 50} (2017), 073001,
\S 8.5.17, eq.\ (8.237); see also \cite{noumiyamada1998painleve,
ohyama2006studies, jimbomiwa1981monodromy}) writes a different
Hamiltonian
\[
  H_{D_6}^{\mathrm{KNY}} \;=\; \tfrac{1}{t}\bigl\{p(p-1) q^{2}
       + (a_1 + a_2) q p + t p - a_2 q\bigr\}
\]
in three root parameters $(a_0, a_1, a_2)$ with affine
constraint $a_0 + a_1 = 1$. The $\tau$-function theory of
Forrester and Witte ({\em Comm.\ Pure Appl.\ Math.}\
{\bf 55} (2002), 679--727, arXiv:math-ph/0201051) writes the
$P_V$ Hamiltonian (loc.\ cit., \S 2.1, eq.\ (2.1)) in {\em four}
parameters $(v_1, v_2, v_3, v_4)$ subject to the null-sum
constraint
\[
  v_1 + v_2 + v_3 + v_4 \;=\; 0
  \qquad\text{(loc.\ cit., \S 2.1, eq.\ (2.2))}.
\]
Under the $P_V \rightsquigarrow P_{III}$ degeneration of FW
loc.\ cit.\ \S 4.1, the $P_{III}$ Hamiltonian is recovered with
two parameters $(v_1, v_2)$ that match Okamoto's
$(\theta_0, \theta_\infty)$ at the WLOG slice
$\eta_0 = \eta_\infty = 1$ by direct comparison of the FW
loc.\ cit.\ eq.\ (4.1) ODE coefficients $\alpha = -4 v_2,\,
\beta = 4(v_1 + 1),\, \gamma = 4,\, \delta = -4$ with the
Okamoto loc.\ cit.\ eq.\ (0.1) form
$\alpha = -4 \eta_\infty \theta_\infty,\,
\beta = 4 \eta_0 (\theta_0 + 1),\, \gamma = 4 \eta_\infty^{2},\,
\delta = -4 \eta_0^{2}$ at $\eta_0 = \eta_\infty = 1$. None of
the three sources --- Okamoto, KNY, FW --- directly furnishes
the project's 4-tuple
$(\alpha_\infty, \alpha_0, \beta_\infty, \beta_0)$ with all
four entries free and named, although FW's $P_V$ 4-tuple
$(v_1, v_2, v_3, v_4)$ provides the null-sum constraint
displayed above.

\medskip
\noindent\textbf{Definition (project rename, trivial form).}%
\hspace{0.5em}
For v1.3 the 4-tuple
$(\alpha_\infty, \alpha_0, \beta_\infty, \beta_0)$ is declared
in terms of the Okamoto Hamiltonian parameters
$(\eta_\infty, \eta_0, \theta_\infty, \theta_0)$ by the
trivial relabel
\begin{align}
  \alpha_\infty   &\;:=\; \eta_\infty,
                                           \tag{3.5.1a}\\
  \alpha_0        &\;:=\; \eta_0,
                                           \tag{3.5.1b}\\
  \beta_\infty    &\;:=\; \theta_\infty,
                                           \tag{3.5.1c}\\
  \beta_0         &\;:=\; \theta_0.
                                           \tag{3.5.1d}
\end{align}
Equations (3.5.1a)--(3.5.1d) carry no additive shift; the
v1.3 rename is a pure symbol substitution. The Okamoto WLOG
normalisation $\eta_\infty = \eta_0 = 1$ is {\em not} imposed
on the project's 4-tuple, since the $\Vquad$ image
$(1/6, 0, 0, -1/2)$ has $\alpha_\infty = 1/6 \ne 1$ and
$\alpha_0 = 0 \ne 1$: the project records all four of
Okamoto's unnormalised Hamiltonian parameters at the
$\Vquad$ point.

\medskip
\noindent\textbf{The $-1/3$ null-sum offset at the $\Vquad$
parameter point.}\hspace{0.5em}
Under (3.5.1a)--(3.5.1d) the $\Vquad$ image yields the partial
sum
\[
  \alpha_\infty + \alpha_0 + \beta_\infty + \beta_0
  \;=\; \tfrac{1}{6} + 0 + 0 + (-\tfrac{1}{2})
  \;=\; -\tfrac{1}{3},
\]
which differs by $-1/3$ from the FW null-sum constraint
$v_1 + v_2 + v_3 + v_4 = 0$ recalled above (FW \S 2.1
eq.\ (2.2)) when applied to the project 4-tuple. The 058R
canonical-map record \cite{siarc_cc_pipeline_g}\footnote{The
follow-up bridge session
{\rm CC-VQUAD-PIII-NORMALIZATION-MAP-RERE-FIRE} of
2026-05-06 (bridge SHA \texttt{2eb9b28}, \S B.3) records the
$-1/3$ partial sum as anomaly~D2, carried forward unchanged
from the 2026-05-02 partial session.} anchors the $-1/3$ value
at the $\Vquad$ parameter point. Three candidate mechanisms
account for the offset:
\begin{enumerate}[label=(\alph*)]
\item {\em Hamiltonian-expansion residual.}\hspace{0.5em}
The auxiliary Hamiltonian $h$ of FW loc.\ cit.\ Proposition 4.1
reads $h = t H + \tfrac{1}{4} v_1^{2} - \tfrac{1}{2} t$
(loc.\ cit.\ eq.\ (4.3)), so $h$ and the primary $t H$ differ
by a constant-in-$(q, p)$, $t$-linear shift $-\tfrac{1}{2} t$
together with a $v_1$-quadratic constant. Pull-back of this
shift along the $\Vquad$ reduction map of
\cite{siarc_cc_pipeline_g} \S B.3 to the $\Vquad$ parameter
point is a candidate origin for the $-1/3$ offset; the
explicit pull-back has not been derived in the present text.
Mechanism (a) is recorded as a hypothesis with the FW
eq.\ (4.3) anchor.
\item {\em Additive-shift component of the rename.}\hspace{0.5em}
If (3.5.1a)--(3.5.1d) are replaced by the non-trivial form
$\alpha_\infty := \eta_\infty + c_{\alpha,\infty}$,
$\alpha_0 := \eta_0 + c_{\alpha,0}$,
$\beta_\infty := \theta_\infty + c_{\beta,\infty}$,
$\beta_0 := \theta_0 + c_{\beta,0}$ with shift sum
$c_{\alpha,\infty} + c_{\alpha,0} + c_{\beta,\infty}
+ c_{\beta,0} = -1/3$, then the $\Vquad$ image lies on the FW
null-sum locus by construction. The symmetric choice $c_{} =
-1/12$ in each of the four coordinates is the simplest
fall-back; this form has no literature anchor and would record
the rename as a project-internal calibration only.
\item {\em Sakai surface-type artefact.}\hspace{0.5em}
The 4-tuple may project root-system data from the Sakai
$D_6^{(1)}$ surface (\cite{ohyama2006studies,
noumiyamada1998painleve}, and the KNY 2017 \S 8.5.17 surface
dictionary cited above) where the affine simple roots satisfy
a non-zero summation constant. Mechanism (c) is gated to a
separate Sakai-surface analysis and is not pursued here; its
derivation routes through the $D_6^{(1)}$ surface machinery
(rational surface, extended affine Weyl, blow-up structure)
outside the scope of \S 3.5.
\end{enumerate}
The selection between (a), (b), and (c) is recorded as a
follow-up to v1.3 and does not affect the \Cref{ex:cc} /
\Cref{thm:vquad-cc} / \Cref{sec:position}(b) material above:
the project's $\Vquad$ image $(1/6, 0, 0, -1/2)$ in
$(\alpha_\infty, \alpha_0, \beta_\infty, \beta_0)$ is fixed
data anchored by \cite{iwasaki1991from, ohyama2006studies},
and the rename (3.5.1a)--(3.5.1d) supplies the symbol
correspondence to the Hamiltonian parameters used in the cited
literature.

\begin{remark}[Symbol-collision note for $(\alpha, \beta)$ in
\S 3.5]
\label{rem:alpha-beta-tuples}
Three distinct $(\alpha, \beta)$-tuples occur in the present
section and the literature cited above:
\begin{itemize}
\item The {\em WKB exponents} $(A, \alpha, \beta, \gamma)$ of
\Cref{prop:wkb}, which label the leading trans-series
structure of $\log|\dn|$ at $n \to \infty$ for degree-$2$
$\PCF$ families.
\item The {\em classical $P_{III}$ ODE coefficients}
$(\alpha, \beta, \gamma, \delta)$ of Okamoto loc.\ cit.\
eq.\ (0.1), with explicit form
$\alpha = -4 \eta_\infty \theta_\infty,\,
\beta = 4 \eta_0 (\theta_0 + 1),\,
\gamma = 4 \eta_\infty^{2},\,
\delta = -4 \eta_0^{2}$. This 4-tuple labels the polynomial
coefficients of the second-order $P_{III}$ ODE in the
dependent variable $w$.
\item The {\em project-internal Okamoto-rename 4-tuple}
$(\alpha_\infty, \alpha_0, \beta_\infty, \beta_0)$ of
(3.5.1a)--(3.5.1d) above, which labels the four Okamoto
Hamiltonian parameters carried over to the project's
$\Vquad$ reduction notation.
\end{itemize}
The three tuples occupy disjoint roles. Cross-citation between
them in the present text is always made via the explicit
subscripts $({}_\infty, {}_0)$ for the Okamoto-rename, the WKB
symbol $A$ for the trans-series leading exponent, and the
polynomial-context label $(\gamma, \delta)$ for the
classical-ODE 4-tuple.
\end{remark}
